\section{本章のまとめ}
\label{sec:quantum-basics-summary}

本章では，量子機械学習を理解するために必要な量子計算の基礎を説明した．
量子ビットの純粋状態は，計算基底$\ket{0}$と$\ket{1}$の複素線形結合として表され，
測定結果の確率は確率振幅の絶対値の二乗によって与えられる．
複数量子ビットの状態空間はテンソル積で構成され，積状態として分解できない状態として
量子もつれが存在する．

量子ゲートはユニタリ変換であり，量子回路は複数の量子ゲートを組み合わせて量子状態を変化させる．
量子回路から得られる期待値は有限回の測定から推定され，量子機械学習モデルの出力として利用できる．
また，古典データを量子回路へ入力するためには，基底符号化，振幅符号化，角度符号化などの
データ符号化が必要になる．

以上の概念を組み合わせることで，古典データ$x$と学習パラメータ$\bm{\theta}$を入力とし，
観測量の期待値を出力するパラメータ付き量子回路を構成できる．
次章では，この回路を教師あり学習の回帰モデルとして学習させる方法を扱う．

